\section*{Overview}

Modular arithmetic is not just ``take \(\% M\) often.'' The important part is knowing which algebraic operations still
behave well after reduction and which ones need extra conditions.

\section*{When to Use It}

Use it whenever:

\begin{itemize}
  \item answers are required modulo \(M\),
  \item intermediate values grow too large,
  \item a recurrence, combinatorial count, or matrix product naturally lives in \(\mathbb{Z}/M\mathbb{Z}\).
\end{itemize}

\section*{Core Idea}

Addition, subtraction, and multiplication respect reduction:
\[
(a+b)\bmod M,\quad (a-b)\bmod M,\quad (ab)\bmod M.
\]
Division does \emph{not} exist automatically. Dividing by \(x\) means multiplying by its modular inverse, which only
makes sense when the inverse exists.

\section*{Key Insight}

The algebraic structure matters. If the modulus is prime, every non-zero residue has an inverse. For composite moduli,
only numbers coprime to the modulus are invertible.

\section*{Worked Problem}

\paragraph{Problem.}
Count the number of ways to build strings of length \(n\), where the recurrence doubles or triples previous counts, and
print the answer modulo \(10^9+7\).

\paragraph{Why modular arithmetic fits.}
The recurrence only needs addition and multiplication, both of which are safe modulo \(M\). So the whole DP can stay in
reduced form from the start.

\section*{Correctness Intuition}

Reducing after each safe operation does not change the final residue because reduction is compatible with addition and
multiplication. The only time care is needed is when an expression secretly relies on division or ordering.

\section*{Complexity Analysis}

The arithmetic operations stay \(O(1)\). The benefit is numerical safety, not asymptotic improvement.

\section*{Implementation}

The code gives a small \texttt{ModInt} template with normalization, addition, subtraction, multiplication, and fast
exponentiation.

\section*{Common Pitfalls}

\begin{itemize}
  \item Using negative residues without re-normalizing into \([0, M-1]\).
  \item Writing \texttt{a / b \% MOD} as if ordinary division were valid.
  \item Overflowing \texttt{a * b} before the modulus is applied.
\end{itemize}

\section*{Variants / Extensions}

\begin{itemize}
  \item Modular inverse and division.
  \item CRT when one modulus is not enough.
  \item NTT-friendly moduli for fast polynomial multiplication.
\end{itemize}

\section*{Practice Problems}

\begin{itemize}
  \item DP or combinatorics modulo a prime.
  \item Fast exponentiation and modular recurrence simulation.
  \item Any counting problem whose raw numbers are too large to store directly.
\end{itemize}

\section*{References}

\begin{itemize}
  \item \href{https://cp-algorithms.com/algebra/module-inverse.html}{cp-algorithms: Modular Inverse and Modular Arithmetic}.
  \item \href{https://atcoder.github.io/ac-library/production/document_en/modint.html}{AtCoder Library: modint}.
  \item \href{https://codeforces.com/catalog}{Codeforces Catalog}.
\end{itemize}
